\paperfigure{figure3_artifact_controls}{fig:p1f3}{Technical and immune-transcript exclusion controls in the discovery cohort. (A) Dissociation standardized mean difference (SMD, 0.246) and the high-state/low-state PTPRC-positive rate ratio (1.60) are divided by rejection bounds 0.8 and 2.0; the dashed line is one. The rounded state cut is 0.184. (B) Each point is one of 25 samples after restricting to 14,561 fibroblasts eligible for panel-count matching to 986 counts per cell. Reprojection preserves sample fractions (ρ=0.984) and 99.5\% of cell labels; the dashed line is identity. (C) Bars are BIC for one Gaussian minus BIC for two Gaussians fitted to sample high-state fractions, using 24 matched samples before and after exclusion. Positive values favour two components. Exclusion retains 11,262 of 19,410 fibroblasts overall, and matched sample fractions correlate at ρ=0.989. This control measures invariance to the exclusion rule; it does not identify every source of ambient RNA or tissue-composition dependence.}
